% !TEX root = ../../main.tex
% File: part2/chapters4/chap4_3.tex

\section{Recipe 3: Fine-tuning một LLM cho Tác vụ Phân loại}
\label{sec:recipe_classification_finetuning}

\textbf{Mục tiêu:} Fine-tune một LLM lớn (ví dụ: Qwen3 8B hoặc Llama 3.1 8B) cho một tác vụ phân loại văn bản một cách hiệu quả về bộ nhớ bằng cách sử dụng QLoRA.

\begin{tcolorbox}[colback=yellow!10!white, colframe=yellow!50!black, title={Khi nào \emph{không} nên dùng công thức này \newtag}, fonttitle=\bfseries]
Với phần lớn bài toán phân loại, một mô hình encoder nhỏ (BERT/RoBERTa/PhoBERT, dưới 1B tham số) được tinh chỉnh toàn bộ vẫn là lựa chọn tốt nhất: rẻ hơn, nhanh hơn, dễ triển khai hơn và thường chính xác không kém. Hãy dùng LLM lớn khi nhãn phụ thuộc vào suy luận phức tạp, khi dữ liệu gán nhãn rất ít (tận dụng kiến thức nền), hoặc khi bạn cần một mô hình đa nhiệm duy nhất.
\end{tcolorbox}

\textbf{Thành phần chính:}
\begin{itemize}
    \item \textbf{`transformers`:} Sử dụng `AutoModelForSequenceClassification` và `Trainer` API.
    \item \textbf{`datasets`:} Để tải và xử lý bộ dữ liệu.
    \item \textbf{`peft`:} Để áp dụng cấu hình LoRA.
    \item \textbf{`bitsandbytes`:} Để lượng tử hóa mô hình sang 4-bit.
\end{itemize}

\textbf{Các bước thực hiện:}
\begin{enumerate}
    \item \textbf{Tải dữ liệu và tokenizer:} Tải một bộ dữ liệu phân loại và tokenizer tương ứng.
    \item \textbf{Tiền xử lý dữ liệu:} Tokenize văn bản.
    \item \textbf{Tải mô hình cơ sở đã được lượng tử hóa:} Tải LLM với `BitsAndBytesConfig` ở chế độ NF4 4-bit.
    \item \textbf{Áp dụng cấu hình PEFT (LoRA):} Bọc mô hình cơ sở bằng một `PeftModel`.
    \item \textbf{Thiết lập `Trainer`:} Định nghĩa `TrainingArguments` và `Trainer`, truyền vào hàm `compute\_metrics`.
    \item \textbf{Huấn luyện:} Bắt đầu quá trình fine-tuning.
\end{enumerate}

\textbf{Mã nguồn:}

\begin{example}{Fine-tuning một LLM cho Tác vụ Phân loại}{ex:finetune_classificationclassification}
    \begin{minted}{python}
    # Bước 1: Cài đặt
    # pip install transformers datasets accelerate evaluate peft bitsandbytes
    import torch
    from datasets import load_dataset
    from transformers import (
        AutoModelForSequenceClassification,
        AutoTokenizer,
        TrainingArguments,
        Trainer,
        BitsAndBytesConfig
    )
    from peft import LoraConfig, get_peft_model
    import numpy as np
    import evaluate
    
    # Bước 2: Tải dữ liệu và tokenizer
    dataset = load_dataset("yelp_review_full")
    model_id = "Qwen/Qwen3-8B"  # hoặc meta-llama/Llama-3.1-8B
    tokenizer = AutoTokenizer.from_pretrained(model_id)
    tokenizer.pad_token = tokenizer.eos_token # nhiều LLM không có pad token riêng
    
    def tokenize_function(examples):
        return tokenizer(examples["text"], padding="max_length", truncation=True)
    
    tokenized_datasets = dataset.map(tokenize_function, batched=True)
    # Giảm kích thước để chạy demo nhanh
    small_train_dataset = tokenized_datasets["train"].shuffle(seed=42).select(range(1000))
    small_eval_dataset = tokenized_datasets["test"].shuffle(seed=42).select(range(1000))
    
    # Bước 3: Tải mô hình cơ sở đã lượng tử hóa
    bnb_config = BitsAndBytesConfig(
        load_in_4bit=True,
        bnb_4bit_quant_type="nf4",
        bnb_4bit_use_double_quant=True,
        bnb_4bit_compute_dtype=torch.bfloat16,
    )
    model = AutoModelForSequenceClassification.from_pretrained(
        model_id,
        num_labels=5, # Yelp review có 5 sao
        quantization_config=bnb_config,
        device_map="auto"
    )
    model.config.pad_token_id = model.config.eos_token_id
    
    # Bước 4: Áp dụng cấu hình PEFT (LoRA)
    lora_config = LoraConfig(r=16, lora_alpha=32, lora_dropout=0.05,
                             target_modules="all-linear", task_type="SEQ_CLS")
    peft_model = get_peft_model(model, lora_config)
    peft_model.print_trainable_parameters()
    
    # Bước 5: Thiết lập Trainer
    metric = evaluate.load("accuracy")
    def compute_metrics(eval_preds):
        logits, labels = eval_preds
        predictions = np.argmax(logits, axis=-1)
        return metric.compute(predictions=predictions, references=labels)
    
    training_args = TrainingArguments(
        output_dir="qlora-classification",
        learning_rate=2e-4,
        per_device_train_batch_size=4,
        per_device_eval_batch_size=4,
        num_train_epochs=1,
        eval_strategy="epoch",
        save_strategy="epoch",
        load_best_model_at_end=True,
    )
    
    trainer = Trainer(
        model=peft_model,
        args=training_args,
        train_dataset=small_train_dataset,
        eval_dataset=small_eval_dataset,
        compute_metrics=compute_metrics,
        processing_class=tokenizer,
    )
    
    # Bước 6: Huấn luyện
    trainer.train()
    \end{minted}
\end{example}
\textbf{Kết quả mong đợi:} Một mô hình LLM lớn được fine-tune hiệu quả trên một GPU tiêu dùng. Chỉ có các trọng số LoRA (vài MB) được lưu lại, thay vì toàn bộ mô hình (vài chục GB).
